
% VLDB template version of 2020-08-03 enhances the ACM template, version 1.7.0:
% https://www.acm.org/publications/proceedings-template
% The ACM Latex guide provides further information about the ACM template

\documentclass[sigconf, nonacm]{acmart}
\usepackage[utf8]{inputenc}
%\documentclass[sigconf, anonymous]{acmart} % replace this with the one above for the camera-ready
\usepackage[utf8]{inputenc}
\DeclareUnicodeCharacter{202F}{\,}



\usepackage{algorithm}
\usepackage{algpseudocode}
\usepackage{subcaption}
\usepackage{placeins}
\usepackage{diagbox}
\usepackage{array}
\usepackage[table]{xcolor}  % For coloring
\usepackage{booktabs}
\usepackage{caption}
\usepackage[utf8]{inputenc}
\usepackage{tikz}
\usepackage[symbol]{footmisc}
\usepackage{float}
\usepackage{needspace}
\usepackage{enumitem}
\usepackage{subcaption}
\usetikzlibrary{shapes.geometric, arrows.meta, positioning, calc, patterns}

\setlength{\abovecaptionskip}{3pt}
\setlength{\belowcaptionskip}{3pt}

%% The following content must be adapted for the final version
% paper-specific
\newcommand\vldbdoi{XX.XX/XXX.XX}
\newcommand\vldbpages{XXX-XXX}
% issue-specific
\newcommand\vldbvolume{14}
\newcommand\vldbissue{1}
\newcommand\vldbyear{2020}
% should be fine as it is
\newcommand\vldbauthors{\authors}
\newcommand\vldbtitle{\shorttitle} 
% leave empty if no availability url should be set
\newcommand\vldbavailabilityurl{URL_TO_YOUR_ARTIFACTS}
% whether page numbers should be shown or not, use 'plain' for review versions, 'empty' for camera ready
\newcommand\vldbpagestyle{plain} 

\settopmatter{authorsperrow=4}

\begin{document}
\title{QVCache: A Query-Aware Vector Cache}

%%
%% The "author" command and its associated commands are used to define the authors and their affiliations.
\author{Anıl Eren Göçer}
\affiliation{%
  \institution{ETH Zurich}
  \city{Zurich}
  \country{Switzerland}
}
\email{agoecer@ethz.ch}


\author{Vasilis Mageirakos}
\affiliation{%
  \institution{ETH Zurich}
  \city{Zurich}
  \country{Switzerland}
}
\email{vmageirakos@inf.ethz.ch}

\author{Yannis Chronis}
\affiliation{%
  \institution{ETH Zurich}
  \city{Zurich}
  \country{Switzerland}
}
\email{chronis@inf.ethz.ch}

\author{Kyoungmin Kim}
\affiliation{%
  \institution{KAIST}
  \city{Daejeon}
  \country{South Korea}
}
\email{kymkim@kaist.ac.kr}

%%
%% The abstract is a short summary of the work to be presented in the
%% article.

 \begin{abstract}
Approximate caching offers a promising way to avoid repeated I/O and computation by serving vector search queries using results from past queries. Existing systems, however, have two main limitations. First, systems such as Proximity and Potluck use a single similarity threshold for all queries, ignoring local neighbor distributions and thus failing to simultaneously achieve high hit rate and recall. Second, all existing systems, including the state-of-the-art vector cache Aker, store each past query together with its top-$k$ results and reuse these results unchanged for similar queries. This frozen-answer design reduces recall and prevents cache hits from composing neighbors across multiple past queries, limiting both recall and hit-rate opportunities. Moreover, these \emph{query-materialized} caches must update cached top-$k$ lists upon insertions, since a newly inserted vector may belong to the top-$k$ results of multiple cached queries. In the worst case, this requires updating all cached entries, resulting in expensive updates or stale results.

We present QVCache, a \emph{data-materialized} approximate cache for vector search that addresses these limitations through two design decisions. First, an online learning algorithm maintains a distinct similarity threshold for each region of the embedding space and dynamically adapts it to the local neighbor distribution. Second, QVCache builds mini-indexes over cached data vectors and searches them to generate answers dynamically, without materializing query vectors or their top-$k$ results. This design enables cache hits to compose neighbors from multiple past queries and allows insertions to be handled efficiently by directly adding new vectors to the mini-indexes.

Integrating QVCache with pgvector, we evaluate it using the benchmark from prior work. QVCache achieves up to $X\%$ higher hit rate and $X\%$ higher recall than the state-of-the-art vector cache, Aker, while reducing average query latency by up to $Z\%$.

\end{abstract}



\maketitle

%%% do not modify the following VLDB block %%
%%% VLDB block start %%%
\pagestyle{\vldbpagestyle}
\begingroup\small\noindent\raggedright\textbf{PVLDB Reference Format:}\\
\vldbauthors. \vldbtitle. PVLDB, \vldbvolume(\vldbissue): \vldbpages, \vldbyear.\\
\href{https://doi.org/\vldbdoi}{doi:\vldbdoi}
\endgroup
\begingroup
\renewcommand\thefootnote{}\footnote{\noindent
This work is licensed under the Creative Commons BY-NC-ND 4.0 International License. Visit \url{https://creativecommons.org/licenses/by-nc-nd/4.0/} to view a copy of this license. For any use beyond those covered by this license, obtain permission by emailing \href{mailto:info@vldb.org}{info@vldb.org}. Copyright is held by the owner/author(s). Publication rights licensed to the VLDB Endowment. \\
\raggedright Proceedings of the VLDB Endowment, Vol. \vldbvolume, No. \vldbissue\ %
ISSN 2150-8097. \\
\href{https://doi.org/\vldbdoi}{doi:\vldbdoi} \\
}\addtocounter{footnote}{-1}\endgroup
%%% VLDB block end %%%

%%% do not modify the following VLDB block %%
%%% VLDB block start %%%
\ifdefempty{\vldbavailabilityurl}{}{
\vspace{.3cm}
\begingroup\small\noindent\raggedright\textbf{PVLDB Artifact Availability:}\\
The source code, data, and/or other artifacts have been made available at \url{\vldbavailabilityurl}.
\endgroup
}
%%% VLDB block end %%%

\section{Introduction}

Modern data-intensive services increasingly rely on Approximate Nearest Neighbor (ANN) search over high-dimensional embeddings. However, scaling ANN search to high recall under strict latency SLOs remains fundamentally constrained by memory capacity and I/O bandwidth.
This tension is now a dominant systems bottleneck: improving recall predictably inflates latency and operational cost, while controlling latency requires either aggressive approximation or prohibitively large memory footprints. Consequently, ANN search has become a critical performance and cost bottleneck in production systems, underpinning recommendation pipelines~\cite{Covington:2016:DNN:YouTube}, search engines~\cite{Liu:2021:Q2S:Facebook, amazon-shop-the-look, visrel}, and Retrieval-Augmented Generation (RAG) workloads for large language models (LLMs)~\cite{Lewis:2020:RAG:NeurIPS, gao2024retrievalaugmentedgenerationlargelanguage}, more recently including agentic applications~\cite{zhao2026amabench, lee2026characterizing}.

This bottleneck is structural rather than incidental. Achieving high recall in ANN search requires traversing increasingly large and irregular neighborhoods in high-dimensional space, leading to random access amplification that cannot be efficiently prefetched or batched \cite{aker}. In graph-based indexes, recall improvements translate directly into expanded graph exploration, candidate explosion, and deeper disk or memory accesses. Although in-memory ANN systems achieve low latency, their memory cost grows linearly with the size of the dataset. Disk-based systems reduce memory pressure, but suffer severe latency degradation at high recall, especially on very large datasets such as billion-scale collections, where pointer chasing and random I/O \cite{jaiswal2022ooddiskannefficientscalablegraph} dominate execution time.  As a result, operators are forced into an undesirable choice between cost and performance, with no mechanism to amortize repeated query work over time. 

Applications that rely on vector search for retrieval exhibit substantial locality. However, unlike the temporal locality observed in traditional database workloads~\cite{chandramouli2018faster}, where the same data objects are repeatedly accessed over time, locality in vector search does not primarily arise from exact query repetition. Instead, repetition occurs in the latent space, where semantically related queries map to nearby embedding regions and retrieve overlapping sets of vectors. We refer to this behavior as \emph{temporal-semantic locality}. A recent study~\cite{lee2026characterizing} reports that vector search queries target overlapping regions of the embedding space and repeatedly retrieve the same vectors, with 73\% of the data vectors accessed by coding agents being accessed repeatedly within and across sessions. This phenomenon also occurs in traditional workloads: up to 40\% and 68\% of user queries in web search and recommender systems, respectively, are semantically repetitive variants of previous queries and thus have nearby vector embeddings~\cite{10.1145/1277741.1277770yahoo, gill2025advancingsemanticcachingllms, li2024recommenderpurposerepeatexploration}.

Systems exploit this locality to reduce query latency through caching. One approach, \emph{traversal-centric caching}, keeps auxiliary index structures in memory, such as data pages, IVF clusters, or graph nodes near the entry point~\cite{pgvector, diskann, 3495724.3496619, 10.1145/3318464.3386131}. Orthogonal to this approach, \emph{result-centric caching} leverages approximate, or similarity, caching~\cite{similarity-caching, similarity-caching-theory-algorithms}. Such systems ~\cite{aker} serve queries directly from cached results of previous queries when these results can answer a new query with sufficient confidence, typically determined by a distance threshold. By bypassing the underlying index traversal, result-centric caching can reduce cache-hit latency by orders of magnitude while introducing negligible memory overhead at the caching layer.


\section{Background} 

\section{Solution Overview}

\section{Algorithms and Operations}



\section{Evaluation Framework}

\section{Experiments}


\section{Related Work}


\section{Conclusion}










\bibliographystyle{ACM-Reference-Format}
\bibliography{references}

\end{document}
\endinput
